\section{Une nouvelle dynamique entre l'Église et le pouvoir impérial}

\subsection*{Introduction}
    Au coeur de la réflexion politique d’Ambroise se présente une complexité intéressante : comment affirmer l'existence d’une autorité de l’Eglise chrétienne, qui soit à la fois complètement autonome vis-à-vis du pouvoir impérial et capable de jouer un rôle politique important dans la gestion de la société, sans pour autant s'emparer d’un pouvoir temporel qui entrerait en conflit avec les prérogatives de l’empereur et de son administration. Si la séparation entre les domaines du politique et du religieux existe depuis le début du christianisme, Ambroise apporte des précisions sur les sphères d’influence, faisant de l’indépendance épiscopale une nécessité pour lancer des actions publiques. L’évêque de Milan réfléchit alors à l’ensemble des relations existantes entre les institutions de l’Eglise et du pouvoir impérial pour développer et légitimer l’autorité de l’Eglise dans la société. Ces relations sont notamment mises en évidence par le corpus de correspondance entre Ambroise et les empereurs. C’est ce qui amène Giuseppe Visonà à qualifier le livre dix des lettres de « manuel de nouvelle doctrine des rapports entre l’Eglise et l’Etat par son principal théoricien. »\footnote{« \latin{E in definitiva un prontuario con la nuova dottrina Chiesa-Stato a firma del suo maggiore elabitore}.» \cite[33]{pizzolato_nec_timeo}.}

\subsection{Introduire une autorité chrétienne}

\subsubsection{\latin{Auctoritas}: Héritage Républicain et appropriation ambrosienne}
Bien que l’on ne puisse pas qualifier de révolution politique l’appropriation du concept d’\latin{auctoritas} par Ambroise et l’Eglise chrétienne, il est tout de même nécessaire de comprendre l’impact qu’a l’utilisation de ce terme par une nouvelle autorité institutionnelle, qui cherche à se détacher des souvenirs glorieux du Sénat de la période républicaine, et de l’actuelle puissance de l’empereur, détenteur officiel de cette fameuse \latin{auctoritas} depuis Auguste. L’évêque de Milan est le premier personnage clérical à utiliser, de façon conséquente, le terme \latin{auctoritas}. Cette notion est par exemple pratiquement absente des traductions de l’Ecriture par Jérôme, alors qu’Ambroise l’utilise pour définir la provenance divine d’une sanction d’ordre spirituel \autocite[214]{sassier_auctoritas}. Pour pleinement saisir la portée de l’appropriation de ce concept par Ambroise, il est nécessaire de revenir certaines structures de la pensée romaine traditionnelle, qui englobe une séparation juridique entre les deux notions d’\latin{auctoritas}, souvent simplement traduit par autorité, et de potestas, pouvant signifier puissance ou pouvoir. L’évêque de Milan est un ancien gouverneur et membre de l’aristocratie romaine, sa pensée politique est donc fondamentalement ancrée dans les traditions romaines. Il connaît parfaitement l’ancienne République et donc l’origine de la notion d’\latin{auctoritas}, ce qui lui permet de l’utiliser avec plus de justesse à son profit.

\bigskip

Il est difficile de comprendre avec précision ce qu’englobe l’\latin{auctoritas} en tant que notion juridique, c’est une idée implantée dans la société romaine qui n’est pas remise en question, mais également jamais franchement définie. Il est par contre clair qu’une séparation importante existe entre l’autorité et le pouvoir, des notions pouvant être dans les mains d’une ou de plusieurs personnes. Les historiens Michel Christol, Frédéric Hurlet et Pierre Cosme en parlent comme ceci : « L’autorité telle que les Romains la concevaient était un surcroît de pouvoir reconnu à un groupe où un individu, conférant une influence et un ascendant à celui qui en était le dépositaire. Elle s’exprimait à travers les initiatives que le groupe ou l’individu prenait et qui étaient suivies d’effet : elle était cette puissance qui permettait de l’emporter dans la prise de décision, sans avoir recours à la force. ». Tout au long de la période républicaine, dans les institutions romaines, l’\latin{auctoritas} est plutôt dans les mains du Sénat, et la potestas, dans celles des magistrats. L’autorité du Sénat n’est donc pas un pouvoir exécutif, mais l’institution peut présenter des décisions « obligatoires », que doivent appliquer les magistrats. C’est bien par leur influence, par l’\latin{auctoritas}, que s’exerce le véritable pouvoir du Sénat. C’est sur cette base juridique que va s’appuyer Ambroise pour instaurer une influence grandissante et pratiquement obligatoire de l’Eglise sur l’empereur.
